%------------------------------------------------%
% Chap. 3
%------------------------------------------------%
\chapter{\bi{Basic linear computation}{基本的な線形計算}}\label{chap:basiclinear}

%------------------------------------------------%
% DD, TD, QD linear computation
%------------------------------------------------%
\section{\bi{Detatypes of D, DD, TD, QD and MPFR vector and matrix}{D, DD, TD, QD および MPFR のベクトルと行列のデータ型}}

\bi{Vector datatypes of DD, TD, QD and MPFR precision are shown as follows:}{DD, TD, QD および MPFR 精度のベクトルのデータ型は次のとおりです。}
\begin{description}
    \item[DDVector, TDVector, QDVector]\strut\par
    \begin{lstlisting}
    typedef struct
    {
        long int dim; // dim <= real_dim
        long int real_dim; // multiplier of _BNC_D_WIDTH
        double *element[DDSIZE]; // [TDSIZE] and [QDSIZE] used as TDVector and QDVector
    } ddvector;

    typedef *DDVector; // DDVector, TDVector, QDVector are pointers.
    \end{lstlisting}
 
    \item[MPFVector]\strut\par
    \begin{lstlisting}
    typedef struct{
        unsigned long int prec; // default precision of element
        mpf_t *element;
        long int dim;
        long int real_dim;
    } mpfvector;
    
    typedef mpfvector *MPFVector; // MPFVector is a pointer.
    \end{lstlisting}
    
\end{description}

\bi{Matrix datatypes of DD, TD, QD and MPFR precision are shown as follows:}{DD, TD, QD および MPFR 精度の行列のデータ型は次のとおりです。}
\begin{description}
    \item[DDMatrix, TDMatrix, QDMatrix]\strut\par
    \begin{lstlisting}
    // DD matrix
    typedef struct{
        long int row_dim, col_dim;
        long int real_row_dim, real_col_dim; // multiplier of _BNC_D_WIDTH
        double *element[DDSIZE];
    } ddmatrix;

    typedef *DDMatrix;
    \end{lstlisting}

    \item[MPFMatrix]\strut\par
    \begin{lstlisting}
    typedef struct{
        unsigned long int prec;
        mpf_t *element;
        long int row_dim, col_dim;
        long int real_row_dim, real_col_dim;
        void *element_block; // mantissa block
    } mpfmatrix;

    typedef mpfmatrix *MPFMatrix;
    \end{lstlisting}
\end{description}

\section{\bi{Vector arithmetic}{ベクトル演算}}

\begin{itemize}
    \funcitem{ddfloat}{get\_ddvector\_i\_ddfloat}{\varnameddvec{vec, long int index}} $\cdots$ \bi{Get index-th element of vec as ddfloat datatype.}{vec の index 番目の要素を ddfloat 型として取得します。}
    \funcitem{tdfloat}{get\_tdvector\_i\_tdfloat}{\varnametdvec{vec, long int index}}
    \funcitem{qdfloat}{get\_qdvector\_i\_qdfloat}{\varnameqdvec{vec, long int index}}

    \macroitem{GET\_DDVECTOR\_I}{vec, index} $\cdots$ \bi{Get index-th element as pointer to ddfloat.}{index 番目の要素を ddfloat へのポインタとして取得します。}
    \macroitem{get\_ddvector\_i}{vec, index}
    \macroitem{GET\_TDVECTOR\_I}{vec, index} $\cdots$ \bi{Get index-th element as pointer to tdfloat.}{index 番目の要素を tdfloat へのポインタとして取得します。}
    \macroitem{get\_tdvector\_i}{vec, index}
    \macroitem{GET\_QDVECTOR\_I}{vec, index} $\cdots$ \bi{Get index-th element as pointer to qdfloat.}{index 番目の要素を qdfloat へのポインタとして取得します。}
    \macroitem{get\_qdvector\_i}{vec, index}
    \macroitem{GET\_MPFVECTOR\_I}{vec, index} $\cdots$ \bi{Get index-th element as pointer to mpf\_t (mpfr\_t).}{index 番目の要素を mpf\_t (mpfr\_t) へのポインタとして取得します。}
    \macroitem{get\_mpfvector\_i}{vec, index}


    \funcitem{void}{set\_ddvector\_i}{\varnameddvec{vec}, \varname{long int}{index}, \varname{double *}{val}} $\cdots$ \bi{Store val to index-th element of vec.}{val を vec の index 番目の要素に格納します。}
    \macroitem{SET\_DDVECTOR\_I}{vec, index, val}
    \funcitem{void}{set\_tdvector\_i}{\varnametdvec{vec}, \varname{long int}{index}, \varname{double *}{val}}
    \macroitem{SET\_TDVECTOR\_I}{vec, index, val}
    \funcitem{void}{set\_qdvector\_i}{\varnameqdvec{vec}, \varname{long int}{index}, \varname{double *}{val}}
    \macroitem{SET\_QDVECTOR\_I}{vec, index, val}
    \funcitem{void}{set\_qdvector\_i}{\varnameqdvec{vec}, \varname{long int}{index}, \varname{double *}{val}}
    \macroitem{SET\_QDVECTOR\_I}{vec, index, val}
    
    \funcitem{void}{set\_ddvector\_i\_d}{\varnameddvec{vec}, \varname{long int}{index}, \varname{double}{val}}$\cdots$ \bi{Store val to index-th element of vec as double.}{val を double として vec の index 番目の要素に格納します。}
    \macroitem{SET\_DDVECTOR\_I\_D}{vec, index, val}
    \funcitem{void}{set\_tdvector\_i\_d}{\varnametdvec{vec}, \varname{long int}{index}, \varname{double}{val}}
    \macroitem{SET\_TDVECTOR\_I\_D}{vec, index, val}
    \funcitem{void}{set\_qdvector\_i\_d}{\varnameqdvec{vec}, \varname{long int}{index}, \varname{double}{val}}
    \macroitem{SET\_QDVECTOR\_I\_D}{vec, index, val}
    
    \funcitem{void}{set0\_ddvector\_i}{\varnameddvec{vec}, \varname{long int}{index}} $\cdots$ \bi{Store zero to index-th element of vec.}{vec の index 番目の要素に 0 を格納します。}
    \macroitem{SET0\_DDVECTOR\_I}{vec, index}
    \funcitem{void}{set0\_tdvector\_i}{\varnametdvec{vec}, \varname{long int}{index}}
    \macroitem{SET0\_TDVECTOR\_I}{vec, index}
    \funcitem{void}{set0\_qdvector\_i}{\varnameqdvec{vec}, \varname{long int}{index}}
    \macroitem{SET0\_QDVECTOR\_I}{vec, index}
    
    
    \funcitem{DDVector}{init\_ddvector}{\varname{in}{dimension}} $\cdots$ \bi{Allocate and initialize a vector and set it as zero vector.}{ベクトルを確保・初期化し、零ベクトルとして設定します。}
    \funcitem{TDVector}{init\_tdvector}{\varname{in}{dimension}}
    \funcitem{QDVector}{init\_qdvector}{\varname{in}{dimension}}
    \funcitem{MPFVector}{init\_mpfvector}{\varname{in}{dimension}} $\cdots$ \bi{Allocate and initialize a MPFVector in default precision in bits with }{}\verb|bnc_set_default_prec|\bi{ and set it as zero vector.}{ によるビット単位の既定精度で MPFVector を確保・初期化し、零ベクトルとして設定します。}
    \funcitem{MPFVector}{init2\_mpfvector}{\varname{int}{dimension}, \varname{unsigned long}{prec}} $\cdots$ \bi{Initialize a MPFVector in prec bits and set it as zero vector.}{prec ビットで MPFVector を初期化し、零ベクトルとして設定します。}
    
    \funcitem{void}{free\_ddvector}{\varnameddvec{vec}} $\cdots$ \bi{Free vec.}{vec を解放します。}
    \funcitem{void}{free\_tdvector}{\varnametdvec{vec}}
    \funcitem{void}{free\_qdvector}{\varnameqdvec{vec}}
    \funcitem{void}{free\_mpfvector}{\varname{MPFVector}{vec}}
    
    \funcitem{void}{set\_ddfloat\_ddvec}{\varname{ddfloat}{ret[]}, \varname{int}{ret\_dim}, \varnameddvec{vec}} $\cdots$ \bi{Convert vec to ddfloat array.}{vec を ddfloat 配列に変換します。}
    \funcitem{void}{set\_tdfloat\_tdvec}{\varname{tdfloat}{ret[]}, \varname{int}{ret\_dim}, \varnametdvec{vec}} $\cdots$ \bi{Convert vec to tdfloat array.}{vec を tdfloat 配列に変換します。}
    \funcitem{void}{set\_qdfloat\_qdvec}{\varname{qdfloat}{ret[]}, \varname{int}{ret\_dim}, \varnameqdvec{vec}} $\cdots$ \bi{Convert vec to qdfloat array.}{vec を qdfloat 配列に変換します。}
    
    \funcitem{void}{set\_ddvector\_ddfloat}{\varnameddvec{ret}, \varname{ddfloat}{array[]}, \varname{int}{array\_dim}} $\cdots$ \bi{Convert ddfloat array to \varnameddvec{ret}.}{ddfloat 配列を \varnameddvec{ret} に変換します。}
    \funcitem{void}{set\_tdvector\_tdfloat}{\varnametdvec{ret}, \varname{tdfloat}{array[]}, \varname{int}{array\_dim}} $\cdots$ \bi{Convert tdfloat array to \varnametdvec{ret}.}{tdfloat 配列を \varnametdvec{ret} に変換します。}
    \funcitem{void}{set\_qdvector\_qdfloat}{\varnameqdvec{ret}, \varname{qdfloat}{array[]}, \varname{int}{array\_dim}} $\cdots$ \bi{Convert qdfloat array to \varnameqdvec{ret}.}{qdfloat 配列を \varnameqdvec{ret} に変換します。}

    \funcitem{void}{print\_ddvector}{\varnameddvec{vec}} $\cdots$ \bi{Print vec.}{vec を表示します。}
    \funcitem{void}{print\_tdvector}{\varnametdvec{vec}}
    \funcitem{void}{print\_qdvector}{\varnameqdvec{vec}}
    \funcitem{void}{print\_mpfvector}{\varname{MPFVector}{vec}}
    
    \funcitem{void}{set0\_ddvector}{\varnameddvec{vec}} $\cdots$ \bi{Set vec as zero vector.}{vec を零ベクトルに設定します。}
    \funcitem{void}{set0\_tdvector}{\varnametdvec{vec}}
    \funcitem{void}{set0\_qdvector}{\varnameqdvec{vec}}
    \funcitem{void}{set0\_mpfvector}{\varname{MPFVector}{vec}}
    
    \funcitem{void}{set\_ddvector\_i\_str}{\varnameddvec{vec}, \varname{long int}{index}, \varname{const char *}{str}} $\cdots$ \bi{Set *str expressed in decimal to index-th element of vec.}{10 進数で表現された *str を vec の index 番目の要素に設定します。}
    \funcitem{void}{set\_tdvector\_i\_str}{\varnametdvec{vec}, \varname{long int}{index}, \varname{const char *}{str}}
    \funcitem{void}{set\_qdvector\_i\_str}{\varnameqdvec{vec}, \varname{long int}{index}, \varname{const char *}{str}}
    \funcitem{void}{set\_mpfvector\_i\_str}{\varnamempfvec{vec}, \varname{long int}{index}, \varname{const char *}{str}}

    \funcitem{void}{add\_ddvector}{\varname{DDVector}{c}, \varname{DDVector}{a}, \varname{DDVector}{b}} $\cdots$  $c := a + b$
    \funcitem{void}{add\_tdvector}{\varname{TDVector}{c}, \varname{TDVector}{a}, \varname{TDVector}{b}}
    \funcitem{void}{add\_qdvector}{\varname{QDVector}{c}, \varname{QDVector}{a}, \varname{QDVector}{b}}
    \funcitem{void}{add\_mpfvector}{\varname{MPFVector}{c}, \varname{MPFVector}{a}, \varname{MPFVector}{b}}

    \funcitem{void}{add2\_ddvector}{\varnameddvec{c}, \varnameddvec{a}} $cdots$ $c := c + a$
    \funcitem{void}{add2\_tdvector}{\varnametdvec{c}, \varnametdvec{a}}
    \funcitem{void}{add2\_qdvector}{\varnametdvec{c}, \varnameqdvec{a}}
    \funcitem{void}{add2\_mpfvector}{\varnamempfvec{c}, \varnamempfvec{a}}

    \funcitem{void}{sub\_ddvector}{\varnameddvec{c}, \varnameddvec{a}, \varnameddvec{b}} $\cdots$ $c := a - b$
    \funcitem{void}{sub\_tdvector}{\varnametdvec{c}, \varnametdvec{a}, \varnametdvec{b}}
    \funcitem{void}{sub\_qdvector}{\varnameqdvec{c}, \varnameqdvec{a}, \varnameqdvec{b}}
    \funcitem{void}{sub\_mpfvector}{\varnamempfvec{c}, \varnamempfvec{a}, \varnamempfvec{b}}

    \funcitem{void}{sub2\_ddvector}{\varnameddvec{c}, \varnameddvec{a}} $\cdots$ $c -= a$
    \funcitem{void}{sub2\_tdvector}{\varnametdvec{c}, \varnametdvec{a}}
    \funcitem{void}{sub2\_qdvector}{\varnameqdvec{c}, \varnameqdvec{a}}
    \funcitem{void}{sub2\_mpfvector}{\varnamempfvec{c}, \varnamempfvec{a}}

    \funcitem{void}{cmul\_ddvector}{\varnameddvec{c}, \varname{double}{val[DDSIZE]}, \varnameddvec{a}} $\cdots$ $c := \mbox{val} \times a$
    \funcitem{void}{cmul\_tdvector}{\varnametdvec{c}, \varname{double}{val[TDSIZE]}, \varnametdvec{a}}
    \funcitem{void}{cmul\_qdvector}{\varnameqdvec{c}, \varname{double}{val[QDSIZE]}, \varnameqdvec{a}}
    \funcitem{void}{cmul\_mpfvector}{\varnamempfvec{c}, \varname{mpf\_t}{val}, \varnamempfvec{a}}

    \funcitem{void}{cmul2\_ddvector}{\varnameddvec{c}, \varname{double}{val[DDSIZE]}} $\cdots$ $c := val \times c$
    \funcitem{void}{cmul2\_ddvector}{\varnametdvec{c}, \varname{double}{val[TDSIZE]}}
    \funcitem{void}{cmul2\_ddvector}{\varnameqdvec{c}, \varname{double}{val[QDSIZE]}}
    \funcitem{void}{cmul2\_ddvector}{\varnamempfvec{c}, \varname{mpf\_t}{val}}

    \funcitem{void}{add\_cmul\_ddvector}{\varnameddvec{c}, \varnameddvec{a}, \varname{double}{val[DDSIZE]}, \varnameddvec{b}} $\cdots$ $c := a + val * b$
    \funcitem{void}{add\_cmul\_tdvector}{\varnametdvec{c}, \varnametdvec{a}, \varname{double}{val[TDSIZE]}, \varnametdvec{b}}
    \funcitem{void}{add\_cmul\_qdvector}{\varnameqdvec{c}, \varnameqdvec{a}, \varname{double}{val[QDSIZE]}, \varnameqdvec{b}} 
    \funcitem{void}{add\_cmul\_mpfvector}{\varnamempfvec{c}, \varnamempfvec{a}, \varname{mpf\_t}{val}, \varnamempfvec{b}} 

    \funcitem{void}{ip\_ddvector}{\varname{double}{ret[DDSIZE]}, \varnameddvec{a}, \varnameddvec{b}} $\cdots$ \bi{Calculate the dot product of $a$ and $b$}{$a$ と $b$ の内積を計算します。}
    \funcitem{void}{ip\_tdvector}{\varname{double}{ret[TDSIZE]}, \varnametdvec{a}, \varnametdvec{b}}
    \funcitem{void}{ip\_qdvector}{\varname{double}{ret[QDSIZE]}, \varnameqdvec{a}, \varnameqdvec{b}}
    \funcitem{void}{ip\_mpfvector}{\varname{mpf\_t}{ret}, \varnamempfvec{a}, \varnamempfvec{b}}
    
    \funcitem{void}{neg\_ddvector}{\varnameddvec{c}, \varnameddvec{a}} $\cdots$  $c := -a$
    \funcitem{void}{neg\_tdvector}{\varnametdvec{c}, \varnametdvec{a}}
    \funcitem{void}{neg\_qdvector}{\varnameqdvec{c}, \varnameqdvec{a}}
    \funcitem{void}{neg\_mpfvector}{\varnamempfvec{c}, \varnamempfvec{a}}

    \funcitem{void}{norm1\_ddvector}{\varname{double}{ret[DDSIZE]}, \varnameddvec{a}} $\cdots$ $\|a\|_1$
    \funcitem{void}{norm1\_tdvector}{\varname{double}{ret[TDSIZE]}, \varnametdvec{a}}
    \funcitem{void}{norm1\_qdvector}{\varname{double}{ret[QDSIZE]}, \varnameqdvec{a}}
    \funcitem{void}{norm1\_mpfvector}{\varname{mpf\_t}{ret}, \varnamempfvec{a}}

    \funcitem{void}{normi\_ddvector}{\varname{double}{ret[DDSIZE]}, \varnameddvec{a}} $\cdots$ $\|a\|_\infty$
    \funcitem{void}{normi\_tdvector}{\varname{double}{ret[TDSIZE]}, \varnametdvec{a}} 
    \funcitem{void}{normi\_qdvector}{\varname{double}{ret[QDSIZE]}, \varnameqdvec{a}} 
    \funcitem{void}{normi\_mpfvector}{\varname{mpf\_t}{ret}, \varnamempfvec{a}} 

    \funcitem{void}{norm2\_ddvector}{\varname{double}{ret[DDSIZE]}, \varnameddvec{vec}} $\cdots$ $\|a\|_2$
    \funcitem{void}{norm2\_tdvector}{\varname{double}{ret[TDSIZE]}, \varnametdvec{vec}}
    \funcitem{void}{norm2\_qdvector}{\varname{double}{ret[QDSIZE]}, \varnameqdvec{vec}}
    \funcitem{void}{norm2\_mpfvector}{\varname{mpf\_t}{ret}, \varnamempfvec{vec}}
    
\end{itemize}

\section{\bi{Matrix arithmetic}{行列演算}}

\begin{itemize}   
    \funcitem{ddfloat}{get\_ddmatrix\_ij\_ddfloat}{\varnameddmat{mat}, \varname{long int}{i}, \varname{long int}{j}} $\cdots$ \bi{Get the $(i, j)$-th element of mat.}{mat の $(i, j)$ 番目の要素を取得します。}
    \macroitem{GET\_DDMATRIX\_IJ}{mat, i, j} ((get\_ddmatrix\_ij\_ddfloat((mat), (i), (j)).val))
    \macroitem{get\_ddmatrix\_ij}{mat, i, j} ((get\_ddmatrix\_ij\_ddfloat((mat), (i), (j)).val))
    \funcitem{tdfloat}{get\_tdmatrix\_ij\_tdfloat}{\varnametdmat{mat}, \varname{long int}{i}, \varname{long int}{j}}
    \macroitem{GET\_TDMATRIX\_IJ}{mat, i, j} ((get\_tdmatrix\_ij\_tdfloat((mat), (i), (j)).val))
    \macroitem{get\_tdmatrix\_ij}{mat, i, j} ((get\_tdmatrix\_ij\_tdfloat((mat), (i), (j)).val))
    \funcitem{qdfloat}{get\_qdmatrix\_ij\_qdfloat}{\varnameqdmat{mat}, \varname{long int}{i}, \varname{long int}{j}}
    \macroitem{GET\_QDMATRIX\_IJ}{mat, i, j} ((get\_qdmatrix\_ij\_qdfloat((mat), (i), (j)).val))
    \macroitem{get\_qdmatrix\_ij}{mat, i, j} ((get\_qdmatrix\_ij\_qdfloat((mat), (i), (j)).val))
    \macroitem{GET\_MPFMATRIX\_IJ}{mat, i, j} ((get\_mpfmatrix\_ij((mat), (i), (j))))
    \funcitem{mpf\_ptr}{get\_mpfmatrix\_ij}{\varnamempfmat{mat}, \varname{long int}{i}, \varname{long int}{j}} 
    
    \funcitem{void}{set\_ddmatrix\_ij}{\varnameddmat{mat}, \varname{long int}{i}, \varname{long int}{j}, \varname{double *}{val}}
    \macroitem{SET\_DDMATRIX\_IJ}{mat, i, j, val} set\_ddmatrix\_ij((mat), (i), (j), (val))
    \funcitem{void}{set\_tdmatrix\_ij}{\varnametdmat{mat}, \varname{long int}{i}, \varname{long int}{j}, \varname{double *}{val}}
    \macroitem{SET\_TDMATRIX\_IJ}{mat, i, j, val} set\_tdmatrix\_ij((mat), (i), (j), (val))
    \funcitem{void}{set\_qdmatrix\_ij}{\varnameqdmat{mat}, \varname{long int}{i}, \varname{long int}{j}, \varname{double *}{val}}
    \macroitem{SET\_QDMATRIX\_IJ}{mat, i, j, val} set\_qdmatrix\_ij((mat), (i), (j), (val))
    \funcitem{void}{set\_mpfmatrix\_ij}{\varnamempfmat{mat}, \varname{long int}{i}, \varname{long int}{j}, \varname{mpf\_t}{val}}
    \macroitem{SET\_MPFMATRIX\_IJ}{mat, i, j, val} set\_mpfmatrix\_ij((mat), (i), (j), (val))
    
    \funcitem{void}{set\_ddmatrix\_ij\_d}{\varnameddmat{mat}, \varname{long int}{i}, \varname{long int}{j}, \varname{double}{val}}
    \macroitem{SET\_DDMATRIX\_IJ\_D}{mat, i, j, val} set\_ddmatrix\_ij\_d((mat), (i), (j), (val))
    \macroitem{SET\_DDMATRIX\_IJ\_UI}{mat, i, j, val} set\_ddmatrix\_ij\_d((mat), (i), (j), (double)(val))
    \macroitem{set\_ddmatrix\_ij\_ui}{mat, i, j, val} set\_ddmatrix\_ij\_d((mat), (i), (j), (double)(val))
    
    \funcitem{void}{set0\_ddmatrix\_ij}{\varnameddmat{mat}, \varname{long int}{i}, \varname{long int}{j}}
    \macroitem{SET0\_DDMATRIX\_IJ}{mat, i, j} set0\_ddmatrix\_ij((mat), (i), (j))
    
    \funcitem{void}{set0\_ddmatrix}{\varnameddmat{mat}} $\cdots$ \bi{set mat as zero matrix}{mat を零行列に設定します。}
    \funcitem{void}{set0\_tdmatrix}{\varnametdmat{mat}}
    \funcitem{void}{set0\_qdmatrix}{\varnameqdmat{mat}}
    \funcitem{void}{set0\_mpfmatrix}{\varnamempfmat{mat}}
    
    \funcitem{DDMatrix}{init\_ddmatrix}{\varname{long int}{row\_dim}, \varname{long int}{col\_dim}} $\cdots$ \bi{Initialize DDMatrix.}{DDMatrix を初期化します。}
    \funcitem{TDMatrix}{init\_tdmatrix}{\varname{long int}{row\_dim}, \varname{long int}{col\_dim}} $\cdots$ \bi{Initialize TDMatrix.}{TDMatrix を初期化します。}
    \funcitem{QDMatrix}{init\_qdmatrix}{\varname{long int}{row\_dim}, \varname{long int}{col\_dim}} $\cdots$ \bi{Initialize QDMatrix.}{QDMatrix を初期化します。}
    \funcitem{MPFMatrix}{init\_mpfmatrix}{\varname{long int}{row\_dim}, \varname{long int}{col\_dim}} $\cdots$ \bi{Initialize MPFMatrix.}{MPFMatrix を初期化します。}
    \funcitem{MPFMatrix}{init2\_mpfmatrix}{\varname{long int}{row\_dim}, \varname{long int}{col\_dim}, \varname{unsigned long}{prec}} $\cdots$ \bi{Initialize MPFMatrix as prec-bit precision.}{MPFMatrix を prec ビット精度で初期化します。}

    
    \funcitem{void}{free\_ddmatrix}{\varnameddmat{mat}} $\cdots$ \bi{Free mat}{mat を解放します。}
    \funcitem{void}{free\_tdmatrix}{\varnametdmat{mat}}
    \funcitem{void}{free\_qdmatrix}{\varnameqdmat{mat}}
    \funcitem{void}{free\_mpfmatrix}{\varnamempfmat{mat}} 
    
    \funcitem{void}{print\_ddmatrix}{\varnameddmat{mat}} $\cdots$ \bi{Print mat.}{mat を表示します。}
    \funcitem{void}{print\_tdmatrix}{\varnametdmat{mat}}
    \funcitem{void}{print\_qdmatrix}{\varnameqdmat{mat}}
    \funcitem{void}{print\_mpfmatrix}{\varnamempfmat{mat}}
    
    \funcitem{void}{set\_ddfloat\_ddmat}{\varname{ddfloat}{ret[]}, \varname{int}{ret\_dim}, \varnameddmat{mat}} $\cdots$ \bi{Convert DDMatrix mat to ddfloat array.}{DDMatrix mat を ddfloat 配列に変換します。}
    \funcitem{void}{set\_tdfloat\_tdmat}{\varname{tdfloat}{ret[]}, \varname{int}{ret\_dim}, \varnameddmat{mat}} $\cdots$ \bi{Convert TDMatrix mat to tdfloat array.}{TDMatrix mat を tdfloat 配列に変換します。}
    \funcitem{void}{set\_qdfloat\_qdmat}{\varname{qdfloat}{ret[]}, \varname{int}{ret\_dim}, \varnameddmat{mat}} $\cdots$ \bi{Convert QDMatrix mat to qdfloat array.}{QDMatrix mat を qdfloat 配列に変換します。}

    \funcitem{void}{set\_ddmatrix\_ddfloat}{\varnameddmat{ret}, \varname{ddfloat}{array[]}, \varname{int}{array\_dim}} $\cdots$ \bi{Convert ddfloat array to DDMatrix.}{ddfloat 配列を DDMatrix に変換します。}
    \funcitem{void}{set\_tdmatrix\_tdfloat}{\varnametdmat{ret}, \varname{tdfloat}{array[]}, \varname{int}{array\_dim}} $\cdots$ \bi{Convert tdfloat array to TDMatrix.}{tdfloat 配列を TDMatrix に変換します。}
    \funcitem{void}{set\_qdmatrix\_qdfloat}{\varnameqdmat{ret}, \varname{qdfloat}{array[]}, \varname{int}{array\_dim}} $\cdots$ \bi{Convert qdfloat array to QDMatrix}{qdfloat 配列を QDMatrix に変換します。}
    
    \funcitem{void}{mul\_ddmatrix}{\varnameddmat{ret}, \varnameddmat{a}, \varnameddmat{b}} $\cdots$ \bi{Simple matrix multiplication.}{単純な行列乗算です。}
    \funcitem{void}{mul\_tdmatrix}{\varnametdmat{ret}, \varnametdmat{a}, \varnametdmat{b}} 
    \funcitem{void}{mul\_qdmatrix}{\varnameqdmat{ret}, \varnameqdmat{a}, \varnameqdmat{b}} 
    \funcitem{void}{mul\_mpfmatrix}{\varnamempfmat{ret}, \varnamempfmat{a}, \varnamempfmat{b}} 
    
    \funcitem{void}{normf\_ddmatrix}{\varname{double}{ret[DDSIZE]}, \varnameddmat{mat}} $\cdots$ \bi{Get the value of Frobenius norm of mat.}{mat のフロベニウスノルムの値を取得します。}
    \funcitem{void}{normf\_tdmatrix}{\varname{double}{ret[TDSIZE]}, \varnametdmat{mat}}
    \funcitem{void}{normf\_qdmatrix}{\varname{double}{ret[QDSIZE]}, \varnameqdmat{mat}}
    \funcitem{void}{normf\_mpfmatrix}{\varname{mpf\_t}{ret}, \varnamempfmat{mat}}
    
    \funcitem{void}{print\_normf\_ddmatrix}{\varname{const char *}{str}, \varnameddmat{mat}} $\cdots$ \bi{print the Frobenius norm of mat.}{mat のフロベニウスノルムを表示します。}
    
    \funcitem{void}{normi\_ddmatrix}{\varname{double}{ret[DDSIZE]}, \varnameddmat{mat}} $\cdots$ \bi{Infinity norm of matrix.}{行列の無限大ノルムです。}
    \funcitem{void}{normi\_tdmatrix}{\varname{double}{ret[TDSIZE]}, \varnametdmat{mat}} 
    \funcitem{void}{normi\_qdmatrix}{\varname{double}{ret[QDSIZE]}, \varnameqdmat{mat}} 
    \funcitem{void}{normi\_mpfmatrix}{\varname{mpf\_t}{ret}, \varnametdmat{mat}} 

    \funcitem{void}{norm1\_ddmatrix}{\varname{double}{ret[DDSIZE]}, \varnameddmat{mat}} $\cdots$ \bi{1-norm of matrix.}{行列の 1 ノルムです。}
    \funcitem{void}{norm1\_tdmatrix}{\varname{double}{ret[TDSIZE]}, \varnametdmat{mat}}
    \funcitem{void}{norm1\_qdmatrix}{\varname{double}{ret[QDSIZE]}, \varnameqdmat{mat}}
    \funcitem{void}{norm1\_mpfmatrix}{\varname{mpf\_t}{ret}, \varnamempfmat{mat}}
    
    \funcitem{void}{add\_ddmatrix}{\varnameddmat{c}, \varnameddmat{a}, \varnameddmat{b}} $\cdots $ $c := a + b$
    \funcitem{void}{add\_tdmatrix}{\varnametdmat{c}, \varnametdmat{a}, \varnametdmat{b}} 
    \funcitem{void}{add\_qdmatrix}{\varnameqdmat{c}, \varnameqdmat{a}, \varnameqdmat{b}} 
    \funcitem{void}{add\_mpfmatrix}{\varnamempfmat{c}, \varnamempfmat{a}, \varnamempfmat{b}} 
    
    \funcitem{void}{sub\_ddmatrix}{\varnameddmat{c}, \varnameddmat{a}, \varnameddmat{b}} $\cdots $ $c := a - b$
    \funcitem{void}{sub\_tdmatrix}{\varnametdmat{c}, \varnametdmat{a}, \varnametdmat{b}}
    \funcitem{void}{sub\_qdmatrix}{\varnameqdmat{c}, \varnameqdmat{a}, \varnameqdmat{b}}
    \funcitem{void}{sub\_mpfmatrix}{\varnamempfmat{c}, \varnamempfmat{a}, \varnamempfmat{b}} 
    
    \funcitem{void}{cmul\_ddmatrix}{\varnameddmat{c}, \varname{double}{sc[DDSIZE]}, \varnameddmat{a}} $\cdots $ $c := sc\times a$
    \funcitem{void}{cmul\_tdmatrix}{\varnametdmat{c}, \varname{double}{sc[TDSIZE]}, \varnametdmat{a}}
    \funcitem{void}{cmul\_ddmatrix}{\varnameqdmat{c}, \varname{double}{sc[QDSIZE]}, \varnameqdmat{a}}
    \funcitem{void}{cmul\_mpfmatrix}{\varnamempfmat{c}, \varname{mpf\_t}{sc}, \varnamempfmat{a}}
    
    \funcitem{void}{transpose\_ddmatrix}{\varnameddmat{c}, \varnameddmat{a}} $\cdots$ $c := a^T$
    \funcitem{void}{transpose\_tdmatrix}{\varnametdmat{c}, \varnametdmat{a}}
    \funcitem{void}{transpose\_qdmatrix}{\varnameqdmat{c}, \varnameqdmat{a}}
    \funcitem{void}{transpose\_mpfmatrix}{\varnamempfmat{c}, \varnamempfmat{a}}
    
    \funcitem{void}{subst\_ddmatrix}{\varnameddmat{c}, \varnameddmat{a}} $\cdots$ $c := a$
    \funcitem{void}{subst\_tdmatrix}{\varnametdmat{c}, \varnametdmat{a}}
    \funcitem{void}{subst\_qdmatrix}{\varnameqdmat{c}, \varnameqdmat{a}}
    \funcitem{void}{subst\_mpfmatrix}{\varnamempfmat{c}, \varnamempfmat{a}} 
    
    \funcitem{void}{setI\_ddmatrix}{\varnameddmat{c}} $\cdots$ $c := I$
    \funcitem{void}{setI\_tdmatrix}{\varnametdmat{c}}
    \funcitem{void}{setI\_qdmatrix}{\varnameqdmat{c}}
    \funcitem{void}{setI\_mpfmatrix}{\varnamempfmat{c}} 
    
    \funcitem{void}{mul\_ddmatrix\_ddvec}{\varnameddvec{v}, \varnameddmat{a}, \varnameddvec{vb}} $\cdots$ $v := a * vb$
    \funcitem{void}{mul\_tdmatrix\_tdvec}{\varnametdvec{v}, \varnametdmat{a}, \varnametdvec{vb}}
    \funcitem{void}{mul\_qdmatrix\_qdvec}{\varnameqdvec{v}, \varnameqdmat{a}, \varnameqdvec{vb}}
    \funcitem{void}{mul\_mpfmatrix\_mpfvec}{\varnamempfvec{v}, \varnamempfmat{a}, \varnamempfvec{vb}} 
    
    \funcitem{void}{mul\_ddmatrixt\_ddvec}{\varnameddvec{v}, \varnameddmat{a}, \varnameddvec{vb}} $\cdots$ $v := a^T * vb$
    \funcitem{void}{mul\_tdmatrixt\_tdvec}{\varnametdvec{v}, \varnametdmat{a}, \varnametdvec{vb}}
    \funcitem{void}{mul\_qdmatrixt\_qdvec}{\varnameqdvec{v}, \varnameqdmat{a}, \varnameqdvec{vb}}
    \funcitem{void}{mul\_mpfmatrixt\_mpfvec}{\varnamempfvec{v}, \varnamempfmat{a}, \varnamempfvec{vb}}
    
    \funcitem{void}{inv\_ddmatrix}{\varnameddmat{a}} $\cdots$ $a := a^{-1}$ (only for square matrix)
    \funcitem{void}{inv\_tdmatrix}{\varnametdmat{a}}
    \funcitem{void}{inv\_qdmatrix}{\varnameqdmat{a}}
    \funcitem{void}{inv\_mpfmatrix}{\varnamempfmat{a}}

\end{itemize}

\begin{itemize}

    \funcitem{void}{subst\_mpfvector\_ddvec}{\varnamempfvec{c}, \varnameddvec{a}}$\cdots$ $c := (mpf)a$   

    \funcitem{void}{subst\_ddvector\_mpfvec}{\varnameddvec{c}, \varnamempfvec{MPFVector a}} $\cdots$ $ c := (dd)a$
    
    \funcitem{void}{subst\_mpfmatrix\_ddmat}{\varnamempfmat{c}, \varnameddmat{a}} $\cdots$ $c := (mpf)a$
    
    \funcitem{void}{subst\_ddmatrix\_mpfmat}{\varnameddmat{c}, \varnamempfmat{a}} $\cdots$ $c := (dd)a$
    

    \funcitem{void}{relerr\_ddvector\_mpfvec}{\varname{double}{relerr[DDSIZE]}, \varnameddvec{approx\_vec}, \varnamempfvec{true\_vec,} \varname{int}{norm\_type}} $\cdots$ \bi{Norm relative error of vector.}{ベクトルのノルム相対誤差です。}

    \funcitem{void}{relerr\_element\_ddvector\_mpf}{\varname{double}{max\_relerr[DDSIZE]}, \varname{double}{min\_relerr[DDSIZE]}, \varname{double}{norm\_relerr[DDSIZE]}, \varnameddvec{approx\_vec}, \varnamempfvec{true\_vec}, \varname{int}{norm\_type}} $\cdots$ \bi{E lementwise relative errors of vector.}{ベクトルの要素ごとの相対誤差です。}
    
    /* c := (dd)a */
    \funcitem{void}{subst\_ddvector\_dvec}{\varnameddvec{c}, \varnamedvec{sa}} $\cdots$ $c := (dd)a$
    
    /* c := (d)a */
    \funcitem{void}{subst\_dvector\_ddvec}{\varnamedvec{c}, \varnameddvec{a}} $\cdots$ $c := (double)a$
    
    /* c := (dd)a */
    \funcitem{void}{subst\_ddmatrix\_dmat}{\varnameddmat{c}, \varnamedmat{a}} $\cdots$ $c := (dd)a$
    
    /* c := (d)a */
    \funcitem{void}{subst\_dmatrix\_ddmat}{\varnamedmat{c}, \varnameddmat{a}} $\cdots$ $c := (double)a$
    
    \funcitem{void}{relerr\_ddvector}{\varname{double}{relerr[DDSIZE]}, \varnameddvec{approx\_vec}, \varnameddvec{true\_vec}, \varname{int}{norm\_type}} $\cdots$ \bi{Norm relative error of vector.}{ベクトルのノルム相対誤差です。}
    
    \funcitem{void}{relerr\_element\_ddvector}{\varname{double}{max\_relerr[DDSIZE]}, \varname{double}{min\_relerr[DDSIZE]}, \varname{double}{norm\_relerr[DDSIZE]}, \varnameddvec{approx\_vec}, \varnameddvec{true\_vec}, \varname{int}{norm\_type}} $\cdots$ \bi{E lementwise relative errors of vector.}{ベクトルの要素ごとの相対誤差です。}
    
    \funcitem{void}{row\_swap\_ddmatrix}{\varnameddmat{mat}, \varname{long int}{row\_index0}, \varname{long int}{row\_index1}, \varname{long int}{col\_start}, \varname{long int}{col\_end}} $\cdots$ \bi{Exchange a(row\_index0, col\_start:col\_end) to a(row\_index1, col\_start:col\_end)}{a(row\_index0, col\_start:col\_end) と a(row\_index1, col\_start:col\_end) を入れ替えます。}
\end{itemize}

\section{\bi{File I/O with matrix and vector}{行列とベクトルのファイル入出力}}

\begin{itemize}
    \funcitem{void}{fread\_ddmatrix}{\varname{FILE *}{fp}, \varnameddmat{mat}} $\cdots$ \bi{Read matrix elements from fp and store theme in mat.}{fp から行列の要素を読み込み、mat に格納します。}
    \funcitem{void}{fread\_tdmatrix}{\varname{FILE *}{fp}, \varnametdmat{mat}}
    \funcitem{void}{fread\_qdmatrix}{\varname{FILE *}{fp}, \varnameqdmat{mat}}
    \funcitem{void}{fread\_mpfmatrix}{\varname{FILE *}{fp}, \varnameqdmat{mat}}

    \funcitem{void}{fread\_ddmatrix\_fname}{\varname{const char *}{fname}, \varnameddmat{mat}} $\cdots$ \bi{Read matrix elements from fname and store theme in mat.}{fname から行列の要素を読み込み、mat に格納します。}
    \funcitem{void}{fread\_tdmatrix\_fname}{\varname{const char *}{fname}, \varnametdmat{mat}}
    \funcitem{void}{fread\_qdmatrix\_fname}{\varname{const char *}{fname}, \varnameqdmat{mat}}
    \funcitem{void}{fread\_mpfmatrix\_fname}{\varname{const char *}{fname}, \varnamempfmat{mat}}


    \funcitem{void}{fwrite\_ddmatrix}{\varname{FILE *}{fp}, \varnameddmat{mat}} $\cdots$ \bi{Write matrix elements of mat to fp.}{mat の行列要素を fp に書き込みます。}
    \funcitem{void}{fwrite\_tdmatrix}{\varname{FILE *}{fp}, \varnametdmat{mat}}
    \funcitem{void}{fwrite\_qdmatrix}{\varname{FILE *}{fp}, \varnameqdmat{mat}}
    \funcitem{void}{fwrite\_mpfmatrix}{\varname{FILE *}{fp}, \varnamempfmat{mat}}


    \funcitem{void}{fwrite\_ddmatrix\_fname}{\varname{const char *}{fname}, \varnameddmat{mat}}$\cdots$ \bi{Write matrix elements of mat to fname.}{mat の行列要素を fname に書き込みます。}
    \funcitem{void}{fwrite\_tdmatrix\_fname}{\varname{const char *}{fname}, \varnametdmat{mat}}
    \funcitem{void}{fwrite\_qdmatrix\_fname}{\varname{const char *}{fname}, \varnameqdmat{mat}}
    \funcitem{void}{fwrite\_mpfmatrix\_fname}{\varname{const char *}{fname}, \varnamempfmat{mat}}

    \funcitem{void}{fread\_ddvector}{\varname{FILE *}{fp}, \varnameddvec{vec}}  $\cdots$ \bi{Read vector elements from fp and store theme in vec.}{fp からベクトルの要素を読み込み、vec に格納します。}
    \funcitem{void}{fread\_tdvector}{\varname{FILE *}{fp}, \varnametdvec{vec}}
    \funcitem{void}{fread\_qdvector}{\varname{FILE *}{fp}, \varnameqdvec{vec}}
    \funcitem{void}{fread\_mpfvector}{\varname{FILE *}{fp}, \varnamempfvec{vec}}

    \funcitem{void}{fread\_ddvector\_fname}{\varname{const char *}{fname}, \varnameddvec{vec}}  $\cdots$ \bi{Read vector elements from fname and store theme in vec.}{fname からベクトルの要素を読み込み、vec に格納します。}
    \funcitem{void}{fread\_tdvector\_fname}{\varname{const char *}{fname}, \varnametdvec{vec}}
    \funcitem{void}{fread\_qdvector\_fname}{\varname{const char *}{fname}, \varnameqdvec{vec}}
    \funcitem{void}{fread\_mpfvector\_fname}{\varname{const char *}{fname}, \varnamempfvec{vec}}

    \funcitem{void}{fwrite\_ddvector}{\varname{FILE *}{fp}, \varnameddvec{vec}} $\cdots$ \bi{Write vector elements of vec to fp.}{vec のベクトル要素を fp に書き込みます。}
    \funcitem{void}{fwrite\_tdvector}{\varname{FILE *}{fp}, \varnametdvec{vec}}
    \funcitem{void}{fwrite\_qdvector}{\varname{FILE *}{fp}, \varnameqdvec{vec}}
    \funcitem{void}{fwrite\_mpfvector}{\varname{FILE *}{fp}, \varnamempfvec{vec}}

    \funcitem{void}{fwrite\_ddvector\_fname}{\varname{const char *}{fname}, \varnameddvec{vec}}$\cdots$ \bi{Write vector elements of vec to fname.}{vec のベクトル要素を fname に書き込みます。}
    \funcitem{void}{fwrite\_tdvector\_fname}{\varname{const char *}{fname}, \varnametdvec{vec}}
    \funcitem{void}{fwrite\_qdvector\_fname}{\varname{const char *}{fname}, \varnameqdvec{vec}}
    \funcitem{void}{fwrite\_mpfvector\_fname}{\varname{const char *}{fname}, \varnamempfvec{vec}}

    
    \funcitem{void}{read\_test\_linear\_eq\_dd}{\varnameddmat{A}, \varnameddvec{true\_x}, \varnameddvec{b}, \varname{long int dim}, \varname{const char *}{fname\_A}, \varname{const char *}{fname\_true\_x}, \varname{const char *}{fname\_b}}  $\cdots$ \bi{Read the coefficient matrix A from fname\_A, the true\_x from fname\_true\_x, and the vector b from fname\_b.}{係数行列 A を fname\_A から、true\_x を fname\_true\_x から、ベクトル b を fname\_b から読み込みます。}
    \funcitem{void}{read\_test\_linear\_eq\_td}{\varnametdmat{A}, \varnametdvec{true\_x}, \varnametdvec{b}, \varname{long int dim}, \varname{const char *}{fname\_A}, \varname{const char *}{fname\_true\_x}, \varname{const char *}{fname\_b}}
    \funcitem{void}{read\_test\_linear\_eq\_qd}{\varnameqdmat{A}, \varnameqdvec{true\_x}, \varnameqdvec{b}, \varname{long int dim}, \varname{const char *}{fname\_A}, \varname{const char *}{fname\_true\_x}, \varname{const char *}{fname\_b}}
    \funcitem{void}{read\_test\_linear\_eq\_mpf}{\varnamempfmat{A}, \varnamempfvec{true\_x}, \varnamempfvec{b}, \varname{long int dim}, \varname{const char *}{fname\_A}, \varname{const char *}{fname\_true\_x}, \varname{const char *}{fname\_b}}

\end{itemize}

\section{\bi{Generating test matrices}{テスト行列の生成}}

\bi{The following functions provide various test matrices for benhmark tests.}{以下の関数は、ベンチマークテスト用にさまざまなテスト行列を提供します。}

\begin{itemize}
    \funcitem{void}{hilbert\_ddmatrix}{\varnameddmat{a}, \varname{long int}{dim}} $\cdots$ \bi{Hilbert matrix.}{ヒルベルト行列です。}
    \funcitem{void}{hilbert\_tdmatrix}{\varnametdmat{a}, \varname{long int}{dim}} 
    \funcitem{void}{hilbert\_qdmatrix}{\varnameqdmat{a}, \varname{long int}{dim}} 
    \funcitem{void}{hilbert\_mpfmatrix}{\varnamempfmat{a}, \varname{long int}{dim}} 

    \funcitem{void}{lotkin\_ddmatrix}{\varnameddmat{a}, \varname{long int}{dim}} $\cdots$ \bi{Lotkin matrix.}{Lotkin 行列です。}
    \funcitem{void}{lotkin\_tdmatrix}{\varnametdmat{a}, \varname{long int}{dim}}
    \funcitem{void}{lotkin\_qdmatrix}{\varnameqdmat{a}, \varname{long int}{dim}} 
    \funcitem{void}{lotkin\_mpfmatrix}{\varnamempfmat{a}, \varname{long int}{dim}}

    \funcitem{void}{frank\_ddmatrix}{\varnameddmat{a}, \varname{long int}{dim}} $\cdots$ \bi{Symmetric Frank matrix.}{対称 Frank 行列です。}
    \funcitem{void}{frank\_tdmatrix}{\varnametdmat{a}, \varname{long int}{dim}}
    \funcitem{void}{frank\_qdmatrix}{\varnameqdmat{a}, \varname{long int}{dim}}
    \funcitem{void}{frank\_mpfmatrix}{\varnamempfmat{a}, \varname{long int}{dim}}

    \funcitem{void}{tridiag\_ddmatrix}{\varnameddmat{a}, \varnameddvec{low\_subdiag}, \varnameddvec{diag}, \varnameddvec{up\_subdiag}, \varname{long int}{dim}} $\cdots$ \bi{Triagonal matrix.}{三重対角行列です。}
    \funcitem{void}{tridiag\_tdmatrix}{\varnameddmat{a}, \varnametdvec{low\_subdiag}, \varnametdvec{diag}, \varnametdvec{up\_subdiag}, \varname{long int}{dim}}
    \funcitem{void}{tridiag\_qdmatrix}{\varnameddmat{a}, \varnameqdvec{low\_subdiag}, \varnameqdvec{diag}, \varnameqdvec{up\_subdiag}, \varname{long int}{dim}}
    \funcitem{void}{tridiag\_mpfmatrix}{\varnameddmat{a}, \varnamempfvec{low\_subdiag}, \varnamempfvec{diag}, \varnamempfvec{up\_subdiag}, \varname{long int}{dim}} 
    
    \funcitem{void}{int\_sym\_rand\_ddmatrix}{\varnameddmat{mat}, \varname{long int}{max}, \varname{long int}{seed}, \varname{long int}{ dim}} $\cdots$ \bi{Integer Symmetrix Random Matrix.}{整数対称ランダム行列です。}
    \funcitem{void}{int\_sym\_rand\_tdmatrix}{\varnametdmat{mat}, \varname{long int}{max}, \varname{long int}{seed}, \varname{long int}{ dim}}
    \funcitem{void}{int\_sym\_rand\_qdmatrix}{\varnameqdmat{mat}, \varname{long int}{max}, \varname{long int}{seed}, \varname{long int}{ dim}}
    \funcitem{void}{int\_sym\_rand\_mpfmatrix}{\varnamempfmat{mat}, \varname{long int}{max}, \varname{long int}{seed}, \varname{long int}{ dim}}
    
    \funcitem{void}{int\_unsym\_rand\_ddmatrix}{\varnameddmat{mat}, \varname{long int}{max}, \varname{long int}{seed}, \varname{long int}{ dim}} $\cdots$ \bi{Integer Unsymmetrix Random Matrix.}{整数非対称ランダム行列です。}
    \funcitem{void}{int\_unsym\_rand\_tdmatrix}{\varnametdmat{mat}, \varname{long int}{max}, \varname{long int}{seed}, \varname{long int}{ dim}} $\cdots$ \bi{Integer Unsymmetrix Random Matrix.}{整数非対称ランダム行列です。}
    \funcitem{void}{int\_unsym\_rand\_qdmatrix}{\varnameqdmat{mat}, \varname{long int}{max}, \varname{long int}{seed}, \varname{long int}{ dim}} $\cdots$ \bi{Integer Unsymmetrix Random Matrix.}{整数非対称ランダム行列です。}
    \funcitem{void}{int\_unsym\_rand\_mpfmatrix}{\varnamempfmat{mat}, \varname{long int}{max}, \varname{long int}{seed}, \varname{long int}{ dim}} $\cdots$ \bi{Integer Unsymmetrix Random Matrix.}{整数非対称ランダム行列です。}
    
    \funcitem{void}{diag\_ddmatrix}{\varnameddmat{mat}, \varnameddvec{diag}, \varname{long int}{dim}} $\cdots$ \bi{Real Diagonal Matrix.}{実対角行列です。}
    \funcitem{void}{diag\_tdmatrix}{\varnametdmat{mat}, \varnametdvec{diag}, \varname{long int}{dim}}
    \funcitem{void}{diag\_qdmatrix}{\varnameqdmat{mat}, \varnameqdvec{diag}, \varname{long int}{dim}}
    \funcitem{void}{diag\_mpfmatrix}{\varnamempfmat{mat}, \varnamempfvec{diag}, \varname{long int}{dim}}
    
    \funcitem{void}{toeplitz\_ddmatrix}{\varnameddmat{mat}, \varname{double}{gamma\_param[DDSIZE]}, \varname{long int}{dim}} $\cdots$ \bi{Toeplitz matrix.}{Toeplitz 行列です。}
    \funcitem{void}{toeplitz\_tdmatrix}{\varnametdmat{mat}, \varname{double}{gamma\_param[TDSIZE]}, \varname{long int}{dim}}
    \funcitem{void}{toeplitz\_qdmatrix}{\varnameqdmat{mat}, \varname{double}{gamma\_param[QDSIZE]}, \varname{long int}{dim}}
    \funcitem{void}{toeplitz\_mpfmatrix}{\varnamempfmat{mat}, \varname{mpf\_t}{gamma\_param}, \varname{long int}{dim}}
    
    \funcitem{void}{pascal\_ddmatrix}{\varnameddmat{ret}, \varname{long int}{dim}} $\cdots $ \bi{Pascal matrix.}{Pascal 行列です。}
    \funcitem{void}{pascal\_tdmatrix}{\varnametdmat{ret}, \varname{long int}{dim}}
    \funcitem{void}{pascal\_qdmatrix}{\varnameqdmat{ret}, \varname{long int}{dim}}
    \funcitem{void}{pascal\_mpfmatrix}{\varnamempfmat{ret}, \varname{long int}{dim}}
    
    \funcitem{void}{im\_rand\_ddmatrix}{\varnameddmat{ret}, \varname{unsigned long}{seed}} $\cdots $  $ I -  random $
    \funcitem{void}{im\_rand\_tdmatrix}{\varnametdmat{ret}, \varname{unsigned long}{seed}}
    \funcitem{void}{im\_rand\_qdmatrix}{\varnameqdmat{ret}, \varname{unsigned long}{seed}}
    \funcitem{void}{im\_rand\_mpfmatrix}{\varnamempfmat{ret}, \varname{unsigned long}{seed}}

\end{itemize}


%------------------------------------------------%
% matrix multiplication
%------------------------------------------------%
\section{\bi{Simple, block, and Strassen matrix multiplication and related functions}{単純法・ブロック法・Strassen 法による行列乗算と関連関数}}

%\begin{lstlisting}  
%    #ifndef _BNC\_DEFAULT\_MIN\_DIM\_STRASSEN
%        \macroitem{\_BNC\_DEFAULT\_MIN\_DIM\_STRASSEN 32
%    #endif // _BNC\_DEFAULT\_MIN\_DIM\_STRASSEN
%    
%    #ifndef STRASSEN\_MIN\_DIM
%    \macroitem{STRASSEN\_MIN\_DIM _BNC\_DEFAULT\_MIN\_DIM\_STRASSEN
%    #endif // STRASSEN\_MIN\_DIM
%\end{lstlisting}

\begin{itemize}

    \macroitem{mul\_ddmatrix\_simple}{c, a, b} $\cdots$ \funcname{mul\_ddmatrix}{c, a, b} $\cdots$ \bi{Matrix multiplication based on simple triple-loop way}{単純な三重ループ方式による行列乗算です。}
    \macroitem{mul\_tdmatrix\_simple}{c, a, b} $\cdots$ \funcname{mul\_tdmatrix}{c, a, b}
    \macroitem{mul\_qdmatrix\_simple}{c, a, b} $\cdots$ \funcname{mul\_qdmatrix}{c, a, b}
    \funcitem{void}{mul\_mpfmatrix\_simple}{\varname{MPFMatrix}{c}, \varname{MPFMatrix}{a}, \varname{MPFMatrix b}}

    \funcitem{void}{mul\_dmatrix\_strassen}{\varname{DMatrix}{ret}, \varname{DMatrix}{mat\_a}, \varname{DMatrix}{mat\_b}, \varname{long int}{min\_dim}} $\cdots$ \bi{Matrix multiplication based on Strassen or Winograd algorithms}{Strassen 法または Winograd 法に基づく行列乗算です。}
    \funcitem{void}{mul\_ddmatrix\_strassen}{\varnameddmat{ret}, \varnameddmat{mat\_a}, \varnameddmat{mat\_b}, \varname{long int}{min\_dim}}
    \funcitem{void}{mul\_tdmatrix\_strassen}{\varnametdmat{ret}, \varnametdmat{mat\_a}, \varnametdmat{mat\_b}, \varname{long int}{min\_dim}}
    \funcitem{void}{mul\_qdmatrix\_strassen}{\varnameqdmat{ret}, \varnameqdmat{mat\_a}, \varnameqdmat{mat\_b}, \varname{long int}{min\_dim}}
    \funcitem{void}{mul\_mpfmatrix\_strassen}{\varname{MPFMatrix}{ret}, \varname{MPFMatrix}{mat\_a}, \varname{MPFMatrix}{mat\_b}, \varname{long int}{min\_dim}} 

    \funcitem{void}{mul\_dmatrix\_block}{DMatrix ret, DMatrix mat\_a, DMatrix mat\_b, \varname{long int}{min\_dim}} $\cdots$ \bi{Block matrix multiplicaiton}{ブロック行列乗算です。}
    \funcitem{void}{mul\_ddmatrix\_block}{\varnameddmat{ret}, \varnameddmat{mat\_a}, \varnameddmat{mat\_b}, \varname{long int}{min\_dim}}
    \funcitem{void}{mul\_tdmatrix\_block}{\varnametdmat{ret}, \varnametdmat{mat\_a}, \varnametdmat{mat\_b}, \varname{long int}{min\_dim}}
    \funcitem{void}{mul\_qdmatrix\_block}{\varnameqdmat{ret}, \varnameqdmat{mat\_}a, \varnameqdmat{mat\_b}, \varname{long int}{min\_dim}}
    \funcitem{void}{mul\_mpfmatrix\_block}{\varname{MPFMatrix}{ret}, \varname{MPFMatrix}{mat\_a}, \varname{MPFMatrix}{mat\_b}, \varname{long int}{min\_dim}}

    \funcitem{void}{mul\_dmatrix\_strassen\_even}{\varname{DMatrix}{ret}, \varname{DMatrix}{mat\_a}, \varname{DMatrix}{mat\_b,} \varname{long int}{min\_dim}} $\cdots$ \bi{Strassen's Algorithm (even dimension)}{Strassen 法（偶数次元）です。}
    \funcitem{void}{mul\_ddmatrix\_strassen\_even}{\varnameddmat{ret}, \varnameddmat{mat\_a}, \varnameddmat{mat\_b}, \varname{long int}{min\_dim}}
    \funcitem{void}{mul\_tdmatrix\_strassen\_even}{\varnametdmat{ret}, \varnametdmat{mat\_a}, \varnametdmat{mat\_b}, \varname{long int}{ min\_dim}}
    \funcitem{void}{mul\_qdmatrix\_strassen\_even}{\varnameqdmat{ret}, \varnameqdmat{mat\_a}, \varnameqdmat{mat\_b}, \varname{long int}{ min\_dim}}
    \funcitem{void}{mul\_mpfmatrix\_strassen\_even}{\varname{MPFMatrix}{ret}, \varname{MPFMatrix}{mat\_a}, \varname{MPFMatrix}{mat\_b}, \varname{long int}{min\_dim}}
    

    \funcitem{void}{mul\_dmatrix\_winograd\_even}{\varname{DMatrix}{ret}, \varname{DMatrix}{mat\_a}, \varname{DMatrix}{mat\_b}, \varname{long int}{min\_dim}} $\cdots$ \bi{Winograd Variant of Strassen's Algorithm}{Strassen 法の Winograd 変種です。}
    \funcitem{void}{mul\_ddmatrix\_winograd\_even}{\varnameddmat{ret}, \varnameddmat{mat\_a}, \varnameddmat{mat\_b}, \varname{long int}{ min\_dim}}
    \funcitem{void}{mul\_tdmatrix\_winograd\_even}{\varnametdmat{ret}, \varnametdmat{mat\_a}, \varnametdmat{mat\_b}, \varname{long int}{min\_dim}}
    \funcitem{void}{mul\_qdmatrix\_winograd\_even}{\varnameqdmat{ret}, \varnameqdmat{mat\_a}, \varnameqdmat{mat\_b}, \varname{long int}{ min\_dim}}     
    \funcitem{void}{mul\_mpfmatrix\_winograd\_even}{\varname{MPFMatrix}{ret}, \varname{MPFMatrix}{mat\_a}, \varname{MPFMatrix}{mat\_b}, \varname{long int}{min\_dim}}
    
    \funcitem{void}{inv\_ddmatrix\_strassen\_even}{\varnameddmat{ret}, \varnameddmat{mat\_a}, \varname{long int}{min\_dim}} $\cdots$ \bi{Computattion of Inverse Matrix by using Strassen's Algorithm}{Strassen 法を用いた逆行列の計算です。}
    \funcitem{void}{inv\_tdmatrix\_strassen\_even}{\varnametdmat{ret}, \varnametdmat{mat\_a}, \varname{long int}{min\_dim}}
    \funcitem{void}{inv\_qdmatrix\_strassen\_even}{\varnameqdmat{ret}, \varnameqdmat{mat\_a}, \varname{long int}{min\_dim}}    

\end{itemize}

\subsection{\bi{OpenMP}{OpenMP}}

\begin{itemize}
   
    \macroitem{\_bncomp\_mul\_ddmatrix\_simple}{c, a, b} \_bncomp\_mul\_ddmatrix(c, a, b) $\cdots$ \bi{Parallelized simple matrix multiplication,}{並列化された単純な行列乗算です。}
    \macroitem{\_bncomp\_mul\_tdmatrix\_simple}{c, a, b} \_bncomp\_mul\_tdmatrix(c, a, b)
    \macroitem{\_bncomp\_mul\_qdmatrix\_simple}{c, a, b} \_bncomp\_mul\_qdmatrix(c, a, b)
    \macroitem{\_bncomp\_mul\_mpfdmatrix\_simple}{c, a, b} \_bncomp\_mul\_mpfdmatrix(c, a, b)
    

    \funcitem{void}{\_bncomp\_mul\_ddmatrix\_block}{\varnameddmat{ret}, \varnameddmat{mat\_a}, \varnameddmat{mat\_b}, \varname{long int}{min\_dim}} $\cdots$ \bi{Parallelized block matrix multiplicaiton.}{並列化されたブロック行列乗算です。}
    \funcitem{void}{\_bncomp\_mul\_tdmatrix\_block}{\varnametdmat{ret}, \varnametdmat{mat\_a}, \varnametdmat{mat\_b}, \varname{long int}{min\_dim}}
    \funcitem{void}{\_bncomp\_mul\_qdmatrix\_block}{\varnameqdmat{ret}, \varnameqdmat{mat\_a}, \varnameqdmat{mat\_b}, \varname{long int}{min\_dim}} 
    \funcitem{void}{\_bncomp\_mul\_mpfmatrix\_block}{\varnamempfmat{ret}, \varnamempfmat{mat\_a}, \varnamempfmat{mat\_b}, \varname{long int}{min\_dim}}


    \funcitem{void}{\_bncomp\_mul\_ddmatrix\_strassen}{\varnameddmat{ret}, \varnameddmat{mat\_a}, \varnameddmat{mat\_b}, \varname{long int}{min\_dim}} $\cdots$ \bi{Parallelized Strassen and Winograd matrix multiplication (limited performance due to poor coding)}{並列化された Strassen 法および Winograd 法による行列乗算です（実装が未成熟なため性能は限定的です）。}
    \funcitem{void}{\_bncomp\_mul\_tdmatrix\_strassen}{\varnametdmat{ret}, \varnametdmat{mat\_a}, \varnametdmat{mat\_b}, \varname{long int}{min\_dim}}
    \funcitem{void}{\_bncomp\_mul\_qdmatrix\_strassen}{\varnameqdmat{ret}, \varnameqdmat{mat\_a}, \varnameqdmat{mat\_b}, \varname{long int}{min\_dim}}
    \funcitem{void}{\_bncomp\_mul\_mpfmatrix\_strassen}{\varnamempfmat{ret}, \varnamempfmat{mat\_a}, \varnamempfmat{mat\_b}, \varname{long int}{min\_dim}}


\end{itemize}

\section{\bi{Getting relative errors of vector and matrix}{ベクトルと行列の相対誤差の取得}}

\begin{itemize}
        \funcitem{void}{relerr3\_dvector}{\varname{double *}{max\_relerr}, \varname{double *}{min\_relerr}, \varname{double *}{norm\_relerr}, \varnamedvec{vec}, \varnamedvec{vec\_true}, \varname{int}{kind\_of\_norm}} $\cdots$ \bi{Relative errors of vec.}{vec の相対誤差です。}
        \funcitem{void}{relerr3\_ddvector}{\varname{double}{max\_relerr[DDSIZE]}, \varname{double}{min\_relerr[DDSIZE]}, \varname{double}{norm\_relerr[DDSIZE]}, \varnameddvec{vec}, \varnameddvec{vec\_true}, \varname{int}{kind\_of\_norm}}
        \funcitem{void}{relerr3\_tdvector}{\varname{double}{max\_relerr[TDSIZE]}}, \varname{double}{min\_relerr[TDSIZE]}, \varname{double}{norm\_relerr[TDSIZE], \varnametdvec{vec}, \varnametdvec{vec\_true}, \varname{int}{kind\_of\_norm}}
        \funcitem{void}{relerr3\_qdvector}{\varname{double}{max\_relerr[QDSIZE]}}, \varname{double}{min\_relerr[QDSIZE]}, \varname{double}{norm\_relerr[QDSIZE], \varnameqdvec{vec}, \varnameqdvec{vec\_true}, \varname{int}{kind\_of\_norm}}
        \funcitem{void}{relerr3\_mpfvector}{\varname{mpf\_t}{max\_relerr}, \varname{mpf\_t}{min\_relerr}, \varname{mpf\_t}{norm\_relerr}, \varnamempfvec{vec}, \varnamempfvec{vec\_true}, \varname{int}{kind\_of\_norm}}

        \funcitem{void}{relerr3\_dmatrix}{\varname{double *}{max\_relerr}, \varname{double *}{min\_relerr}, \varname{double *}{norm\_relerr}, \varname{DMatrix}{mat}, \varname{DMatrix}{mat\_true}, \varname{int}{kind\_of\_norm}} $\cdots$ \bi{Relative errors of mat.}{mat の相対誤差です。}
        \funcitem{void}{relerr3\_ddmatrix}{\varname{double}{max\_relerr[DDSIZE]}, \varname{double}{min\_relerr[DDSIZE]}, \varname{double}{norm\_relerr[DDSIZE]}, \varnameddmat{mat}, \varnameddmat{mat\_true}, \varname{int}{kind\_of\_norm}}
        \funcitem{void}{relerr3\_tdmatrix}{\varname{double}{max\_relerr[TDSIZE]}, \varname{double}{min\_relerr[TDSIZE]}, \varname{double}{norm\_relerr[TDSIZE]}, \varnametdmat{mat}, \varnametdmat{mat\_true}, \varname{int}{kind\_of\_norm}}
        \funcitem{void}{relerr3\_qdmatrix}{\varname{double}{max\_relerr[QDSIZE]}, \varname{double}{min\_relerr[QDSIZE]}, \varname{double}{norm\_relerr[QDSIZE]}, \varnameqdmat{mat}, \varnameqdmat{mat\_true}, \varname{int}{kind\_of\_norm}}
        \funcitem{void}{relerr3\_mpfmatrix}{\varname{mpf\_t}{max\_relerr}, \varname{mpf\_t}{min\_relerr}, \varname{mpf\_t}{norm\_relerr}, \varname{MPFMatrix}{mat}, \varname{MPFMatrix}{mat\_true}, \varname{int}{kind\_of\_norm}}
\end{itemize}

%------------------------------------------------%
% matrix multiplication
%------------------------------------------------%
\section{\bi{Ozaki scheme and related functions}{尾崎スキームと関連関数}}

\bi{These following functions are necessary for Ozaki scheme and multi-fold precision arithmetic.}{以下の関数は、尾崎スキームおよび多倍長精度演算に必要です。}

\begin{itemize}

    \funcitem{void}{extract\_dvector}{\varname{DVector}{ret\_high\_vec}, \varname{DVector}{ret\_low\_vec}, \varname{DVector}{org\_vec}, \varname{double}{num\_bits}}$\cdots$ \bi{ret\_high + ret\_low = org\_vec}{ret\_high + ret\_low = org\_vec となります。}
    \funcitem{void}{split\_dmatrix}{\varname{DMatrix ret\_high\_mat}, \varname{DMatrix}{ret\_low\_mat}, \varname{DMatrix org\_mat}}$\cdots$ \bi{split org\_mat to ret\_high\_mat and ret\_low\_mat in row direction}{org\_mat を行方向に ret\_high\_mat と ret\_low\_mat に分割します。}
    \funcitem{void}{split\_dmatrix\_t}{\varname{DMatrix ret\_high\_mat},\varname{DMatrix}{ret\_low\_mat}, \varname{DMatrix org\_mat}}$\cdots$ \bi{split org\_mat to ret\_high\_mat and ret\_low\_mat in column direction}{org\_mat を列方向に ret\_high\_mat と ret\_low\_mat に分割します。}

    \funcitem{int}{split\_ddvector\_dvec}{\varname{DVector}{ret\_vec[]}, \varname{int}{num\_div}, \varname{DDVector}{org\_vec}} $\cdots$
     \bi{Split org\_vec to ret\_vec[0] $+$ ret\_vec[1] $+ ... +$ ret\_vec[num\_div - 1]}{org\_vec を ret\_vec[0] $+$ ret\_vec[1] $+ ... +$ ret\_vec[num\_div - 1] に分割します。}
    \funcitem{int}{split\_tdvector\_dvec}{\varname{DVector}{ret\_vec[]}, \varname{int}{num\_div}, \varname{TDVector}{org\_vec}} 

    \funcitem{void}{absmax\_ddvector}{\varname{double}{ret[DDSIZE]}, \varname{long int *}{max\_index}, \varname{DDVector}{vec}} $\cdots$ \bi{get absolutely maximum of elements in vec.}{vec の要素のうち絶対値が最大のものを取得します。}
    \funcitem{void}{absmax\_tdvector}{\varname{double}{ret[TDSIZE]}, \varname{long int *}{max\_index}, \varname{TDVector}{vec}}
    \funcitem{void}{absmax\_qdvector}{\varname{double}{ret[QDSIZE]}, \varname{long int *}{max\_index}, \varname{QDVector}{vec}}
    \funcitem{void}{absmax\_mpfvector}{\varname{mpf\_t}{ret}, \varname{long int *}{max\_index}, \varname{MPFVector}{vec}}

    \funcitem{void}{add\_ddvector\_dvec}{\varname{DDVector}{c}, \varname{DDVector}{a}, \varname{DVector}{b}} $\cdots$  $ c := a + (double)b$
    \funcitem{void}{add\_tdvector\_dvec}{\varname{TDVector}{c}, \varname{TDVector}{a}, \varname{DVector}{b}}
    \funcitem{void}{add\_qdvector\_dvec}{\varname{QDVector}{c}, \varname{QDVector}{a}, \varname{DVector}{b}}
    \funcitem{void}{add\_mpfvector\_dvec}{\varname{MPFVector}{c}, \varname{MPFVector}{a}, \varname{DVector}{b}};

    \funcitem{void}{sub\_ddvector\_dvec}{\varname{DDVector}{c}, \varname{DDVector}{a}, \varname{DVector}{b}} $\cdots$  $c := a - (double)b$
    \funcitem{void}{sub\_tdvector\_dvec}{\varname{TDVector}{c}, \varname{TDVector}{a}, \varname{DVector}{b}}
    \funcitem{void}{sub\_qdvector\_dvec}{\varname{QDVector}{c}, \varname{QDVector}{a}, \varname{DVector}{b}}
    \funcitem{void}{sub\_mpfvector\_dvec}{\varname{MPFVector}{c}, \varname{MPFVector}{a}, \varname{DVector}{b}}

    \funcitem{void}{subst\_dvector\_qdvec}{\varname{DVector}{c}, \varname{QDVector}{a}} $\cdots$ $c := (DDVector)a$

    \funcitem{void}{absmax\_row\_ddmatrix}{\varname{double}{mu[DDSIZE]}, \varname{long int *}{max\_j}, \varname{long int}{row\_index}, \varname{DDMatrix}{mat}} $\cdots$ \bi{Return the absolutely maximum element in the max\_j-th row.}{max\_j 番目の行で絶対値が最大の要素を返します。}
    \funcitem{void}{absmax\_row\_tdmatrix}{\varname{double}{mu[TDSIZE]}, \varname{long int *}{max\_j}, \varname{long int}{row\_index}, \varname{TDMatrix}{mat}}
    \funcitem{void}{absmax\_row\_qdmatrix}{\varname{double}{mu[QDSIZE]}, \varname{long int *}{max\_j}, \varname{long int}{row\_index}, \varname{QDMatrix}{mat}}
    \funcitem{void}{absmax\_row\_mpfmatrix}{\varname{mpf\_t}{mu}, \varname{long int *}{max\_j}, \varname{long int}{row\_index}, \varname{MPFMatrix}{mat}}


    \funcitem{void}{add\_ddmatrix\_dmat}{\varname{DDMatrix}{c}, \varname{DDMatrix}{a}, \varname{DMatrix}{b}} $\cdots$  $c := a + (DMatrix)b$
    \funcitem{void}{add\_tdmatrix\_dmat}{\varname{TDMatrix}{c}, \varname{TDMatrix}{a}, \varname{DMatrix}{b}}
    \funcitem{void}{add\_qdmatrix\_dmat}{\varname{QDMatrix}{c}, \varname{QDMatrix}{a}, \varname{DMatrix}{b}}
    \funcitem{void}{add\_mpfmatrix\_dmat}{\varname{MPFMatrix}{c}, \varname{MPFMatrix}{a}, \varname{DMatrix}{b}}

    \funcitem{void}{sub\_ddmatrix\_dmat}{\varname{DDMatrix}{c}, \varname{DDMatrix}{a}, \varname{DMatrix}{b}} $\cdots$  $c := a - (DMatrix)b$
    \funcitem{void}{sub\_tdmatrix\_dmat}{\varname{TDMatrix}{c}, \varname{TDMatrix}{a}, \varname{DMatrix}{b}}
    \funcitem{void}{sub\_qdmatrix\_dmat}{\varname{QDMatrix}{c}, \varname{QDMatrix}{a}, \varname{DMatrix}{b}}
    \funcitem{void}{sub\_mpfmatrix\_dmat}{\varname{MPFMatrix}{c}, \varname{MPFMatrix}{a}, \varname{DMatrix}{b}}

    \funcitem{int}{split\_ddmatrix\_dmat}{\varname{DMatrix}{ret\_mat[]}, \varname{int}{num\_div}, \varname{DDMatrix}{org\_mat}} $cdots$ \bi{Split org\_mat to ret\_mat[] in row direction and return real number of divisions}{org\_mat を行方向に ret\_mat[] に分割し、実際の分割数を返します。}
    \funcitem{int}{split\_tdmatrix\_dmat}{\varname{DMatrix}{ret\_mat[]}, \varname{int}{num\_div}, \varname{TDMatrix}{org\_mat}}
    \funcitem{int}{split\_qdmatrix\_dmat}{\varname{DMatrix}{ret\_mat[]}, \varname{int}{num\_div}, \varname{QDMatrix}{org\_mat}}
    \funcitem{int}{split\_mpfmatrix\_dmat}{\varname{DMatrix}{ret\_mat[]}, \varname{int}{num\_div}, \varname{MPFMatrix}{org\_mat}}

    \funcitem{void}{absmax\_col\_ddmatrix}{\varname{double}{mu[DDSIZE]}, \varname{long int *}{max\_i}, \varname{long int}{col\_index}, \varname{DDMatrix}{mat}} $\cdots$ \bi{return absolutely maximum element and its $(i, j)$-index in each column of mat.}{mat の各列について絶対値が最大の要素とその $(i, j)$ インデックスを返します。}
    \funcitem{void}{absmax\_col\_tdmatrix}{\varname{double}{mu[TDSIZE]}, \varname{long int *}{max\_i}, \varname{long int}{col\_index}, \varname{TDMatrix}{mat}}
    \funcitem{void}{absmax\_col\_qdmatrix}{\varname{double}{mu[QDSIZE]}, \varname{long int *}{max\_i}, \varname{long int}{col\_index}, \varname{QDMatrix{mat}}}
    \funcitem{void}{absmax\_col\_mpfmatrix}{\varname{mpf\_t}{mu}, \varname{long int *}{max\_i}, \varname{long int}{col\_index}, \varname{MPFMatrix}{mat}}

    \funcitem{int}{split\_ddmatrix\_t\_dmat}{\varname{DMatrix}{ret\_mat[]}, \varname{int}{num\_div}, \varname{DDMatrix}{org\_mat}} $\cdots$ \bi{Split org\_mat to ret\_mat[] in column direction and return real number of divisions.}{org\_mat を列方向に ret\_mat[] に分割し、実際の分割数を返します。}
    \funcitem{int}{split\_tdmatrix\_t\_dmat}{\varname{DMatrix}{ret\_mat[]}, \varname{int}{num\_div}, \varname{TDMatrix}{org\_mat}}
    \funcitem{int}{split\_qdvector\_dvec}{\varname{DVector}{ret\_vec[]}, \varname{int}{num\_div}, \varname{QDVector}{org\_vec}}
    \funcitem{int}{split\_qdmatrix\_t\_dmat}{\varname{DMatrix}{ret\_mat[]}, \varname{int}{num\_div}, \varname{QDMatrix}{org\_mat}}

    \funcitem{int}{split\_mpfmatrix\_t\_dmat}{\varname{DMatrix}{ret\_mat[]}, \varname{int num\_div}, \varname{MPFMatrix}{org\_mat}}

    \funcitem{void}{subst\_qdmatrix\_dmat}{\varname{QDMatrix}{c}, \varname{DMatrix}{a}}
    \funcitem{void}{subst\_dmatrix\_qdmat}{\varname{DMatrix}{c}, \varname{QDMatrix}{a}}



    \funcitem{void}{mul\_ddmatrix\_oz}{\varname{DDMatrix}{ret}, \varname{DDMatrix}{a}, \varname{int}{max\_num\_div\_a}, \varname{DDMatrix}{b}, \varname{int}{max\_num\_div\_b}} $cdots$ \bi{Matrix multiplication based on Ozaki scheme}{尾崎スキームに基づく行列乗算です。}
    \funcitem{void}{mul\_tdmatrix\_oz}{\varname{TDMatrix}{ret}, \varname{TDMatrix}{a}, \varname{int}{max\_num\_div\_a}, \varname{TDMatrix}{b}, \varname{int}{max\_num\_div\_b}}
    \funcitem{void}{mul\_qdmatrix\_oz}{\varname{QDMatrix}{ret}, \varname{QDMatrix}{a}, \varname{int}{max\_num\_div\_a}, \varname{QDMatrix}{b}, \varname{int}{max\_num\_div\_b}}
    \funcitem{void}{mul\_mpfmatrix\_oz}{\varname{MPFMatrix}{ret}, \varname{MPFMatrix}{a}, \varname{int}{max\_num\_div\_a}, \varname{MPFMatrix}{b}, \varname{int}{max\_num\_div\_b}}

    \funcitem{void}{mul\_ddmatrix\_ddvec\_oz}{\varname{DDVector}{ret}, \varname{DDMatrix}{a}, \varname{int}{max\_num\_div\_a}, \varname{DDVector} {vb}, \varname{int}{max\_num\_div\_vb}} $\cdots$ \bi{Matrix-Vector multiplication based on Ozaki scheme}{尾崎スキームに基づく行列ベクトル乗算です。}
    \funcitem{void}{mul\_tdmatrix\_tdvec\_oz}{\varname{TDVector}{ret}, \varname{TDMatrix}{a}, \varname{int}{max\_num\_div\_a}, \varname{TDVector}{vb}, \varname{int}{max\_num\_div\_vb}}
    \funcitem{void}{mul\_qdmatrix\_qdvec\_oz}{\varname{QDVector}{ret}, \varname{QDMatrix}{a}, \varname{int}{max\_num\_div\_a}, \varnameqdvec{vb}, \varname{int}{max\_num\_div\_vb}}
    \funcitem{void}{mul\_mpfmatrix\_mpfvec\_oz}{\varname{MPFVector}{ret}, \varname{MPFMatrix}{a}, \varname{int}{max\_num\_div\_a}, \varname{MPFVector}{vb}, \varname{int}{max\_num\_div\_vb}}

    \funcitem{void}{mul\_dmatrix\_oz}{\varname{DMatrix}{ret}, \varname{DMatrix}{a}, \varname{int}{max\_num\_div\_a}, \varname{DMatrix}{b}, \varname{int}{max\_num\_div\_b}} $\cdots$ \bi{The double-precision member of the family; it obeys \texttt{max\_num\_div\_*} exactly as the others do.}{本ファミリの倍精度版です。他と同様に \texttt{max\_num\_div\_*} に従います。}

\end{itemize}

\subsection{\bi{Parallel execution and tuning}{並列実行とチューニング}}

\bi{All the routines above are parallelized with OpenMP. They spend their time in three places --- splitting the operands into double-precision slices, multiplying the slices, and accumulating the slice products back in multi-component or \texttt{mpf\_t} arithmetic --- and all three are parallel. The rows of the result are cut into blocks; one thread takes a block at a time, runs every slice product for it with a \emph{single-threaded} BLAS call and accumulates on the spot. The block loop is dynamically scheduled, which keeps cores of unequal speed busy.}{上記の関数はすべてOpenMPで並列化されています。処理時間は、オペランドを倍精度スライスに分割する部分、スライス同士を乗算する部分、そしてスライス積を多倍長演算で再集約する部分の3箇所に分かれますが、その3つすべてが並列化されています。結果行列の行をブロックに分割し、1スレッドが1ブロックを担当して、そのブロックに対する全スライス積を\emph{シングルスレッドの}BLAS呼び出しで計算し、その場で集約します。ブロックループは動的スケジューリングのため、性能の異なるコアが混在していても待ちが生じません。}

\bi{The row blocks are disjoint and every element of the result still sums its slice products in the original order, so the result does not depend on the number of threads.}{行ブロックは互いに素であり、結果の各要素は従来と同じ順序でスライス積を加算するため、結果はスレッド数に依存しません。}

\bi{Everything is tunable at run time, so a build never has to be repeated just to try another setting:}{すべて実行時に設定できるため、設定を変えるためにビルドし直す必要はありません。}

\begin{itemize}
    \item \texttt{BNC\_OZ\_NUM\_THREADS} $\cdots$ \bi{threads used by the kernels (default: the OpenMP maximum). \texttt{1} disables the kernels' own parallelism and hands the products back to the BLAS.}{関数が使用するスレッド数（既定値：OpenMPの最大値）。\texttt{1} にすると関数自身の並列化を無効にし、乗算をBLASに任せます。}
    \item \texttt{BNC\_OZ\_BLAS\_THREADS} $\cdots$ \bi{threads left to the CBLAS backend while a kernel runs its own parallel loop (default \texttt{1}).}{関数が自前の並列ループを実行している間にCBLASバックエンドへ許すスレッド数（既定値 \texttt{1}）。}
    \item \texttt{BNC\_OZ\_BLOCK\_ROWS} $\cdots$ \bi{rows of the result per block (default: automatic).}{ブロックあたりの結果行数（既定値：自動）。}
    \item \texttt{BNC\_OZ\_BLOCKS\_PER\_THREAD} $\cdots$ \bi{blocks per thread when the block size is automatic (default \texttt{2}).}{ブロックサイズが自動のときのスレッドあたりブロック数（既定値 \texttt{2}）。}
    \item \texttt{BNC\_OZ\_GEMM\_MODE} $\cdots$ \bi{\texttt{own} (parallel block loop with a serial BLAS) or \texttt{blas} (serial block loop with a threaded BLAS); default \texttt{own}.}{\texttt{own}（並列ブロックループ＋直列BLAS）または \texttt{blas}（直列ブロックループ＋並列BLAS）。既定値は \texttt{own} です。}
\end{itemize}

\bi{The same settings are reachable from C as \texttt{bnc\_oz\_set\_num\_threads()}, \texttt{bnc\_oz\_set\_blas\_threads()}, \texttt{bnc\_oz\_set\_block\_rows()} and \texttt{bnc\_oz\_set\_gemm\_mode()}, declared in \texttt{oz\_scheme.h}. Binding the threads with \texttt{OMP\_PROC\_BIND=close OMP\_PLACES=cores} is worth 10--20\%; it is also what makes a BLAS-threaded fallback unusable, because OpenBLAS's pthread pool inherits the affinity mask of the thread that creates it.}{同じ設定はCからも \texttt{bnc\_oz\_set\_num\_threads()}、\texttt{bnc\_oz\_set\_blas\_threads()}、\texttt{bnc\_oz\_set\_block\_rows()}、\texttt{bnc\_oz\_set\_gemm\_mode()} で行えます（\texttt{oz\_scheme.h} で宣言）。\texttt{OMP\_PROC\_BIND=close OMP\_PLACES=cores} によるスレッド固定は10〜20\%の効果があります。一方で、これはBLAS並列に頼るフォールバックを使い物にならなくする要因でもあります。OpenBLASのpthreadプールは生成元スレッドのアフィニティマスクを継承するためです。}

\subsection{\bi{Exponent range}{指数部の範囲}}

\bi{The slices are plain \texttt{double} matrices, but the operands need not fit in the \texttt{double} exponent range: \texttt{mpf\_t} reaches $2^{\pm 2^{30}}$, and even a DD matrix may hold entries near \texttt{DBL\_MAX}. Each row of $A$ and each column of $B$ is therefore scaled by a power of two before it is split, and the exponent is carried along with the slice:}{スライスは通常の \texttt{double} 行列ですが、オペランドが \texttt{double} の指数範囲に収まるとは限りません。\texttt{mpf\_t} は $2^{\pm 2^{30}}$ に達し、DD行列でも \texttt{DBL\_MAX} 近傍の要素を持ち得ます。そこで、分割前に $A$ の各行と $B$ の各列を2の冪でスケーリングし、その指数をスライスとともに保持します。}
\begin{equation*}
C_{ij} = \sum_{p,q} 2^{\,\sigma^A_{p,i} + \sigma^B_{q,j}}\;\bigl(A^{(p)} B^{(q)}\bigr)_{ij}
\end{equation*}
\bi{The exponent is constant along the contracted index, so it factors straight out of the inner product and the error-free property of the slice products is untouched; scaling by a power of two is exact in both \texttt{double} and \texttt{mpf\_t}. A fresh exponent is taken for every slice, so the dynamic range a single row may span is bounded by \texttt{max\_num\_div} alone.}{指数は縮約される添字について定数なので、内積の外に完全に括り出され、スライス積の誤差なし性は損なわれません。2の冪倍は \texttt{double} でも \texttt{mpf\_t} でも厳密です。スライスごとに指数を取り直すため、1つの行が跨げるダイナミックレンジは \texttt{max\_num\_div} のみで決まります。}

\bi{The split functions hand the exponents back through \texttt{\_ex} entry points, each taking a \texttt{long int} array of \texttt{num\_div} $\times$ (row or column count). The names without \texttt{\_ex} remain and forward with a \texttt{NULL} array, which asks for the unscaled split.}{分割関数は \texttt{\_ex} 版で指数を返します。引数は \texttt{num\_div} $\times$（行数または列数）の \texttt{long int} 配列です。\texttt{\_ex} の付かない従来名も残っており、\texttt{NULL} を渡してスケーリングなしの分割を行います。}

\begin{itemize}
    \funcitem{int}{split\_dmatrix\_dmat\_ex}{\varname{DMatrix}{ret\_mat[]}, \varname{long int}{row\_shift[]}, \varname{int}{num\_div}, \varname{DMatrix}{org\_mat}}
    \funcitem{int}{split\_ddmatrix\_dmat\_ex}{\varname{DMatrix}{ret\_mat[]}, \varname{long int}{row\_shift[]}, \varname{int}{num\_div}, \varname{DDMatrix}{org\_mat}}
    \funcitem{int}{split\_tdmatrix\_dmat\_ex}{\varname{DMatrix}{ret\_mat[]}, \varname{long int}{row\_shift[]}, \varname{int}{num\_div}, \varname{TDMatrix}{org\_mat}}
    \funcitem{int}{split\_qdmatrix\_dmat\_ex}{\varname{DMatrix}{ret\_mat[]}, \varname{long int}{row\_shift[]}, \varname{int}{num\_div}, \varname{QDMatrix}{org\_mat}}
    \funcitem{int}{split\_mpfmatrix\_dmat\_ex}{\varname{DMatrix}{ret\_mat[]}, \varname{long int}{row\_shift[]}, \varname{int}{num\_div}, \varname{MPFMatrix}{org\_mat}}
    \funcitem{int}{split\_dmatrix\_t\_dmat\_ex}{\varname{DMatrix}{ret\_mat[]}, \varname{long int}{col\_shift[]}, \varname{int}{num\_div}, \varname{DMatrix}{org\_mat}} $\cdots$ \bi{and the \texttt{dd}, \texttt{td}, \texttt{qd}, \texttt{mpf} counterparts}{および \texttt{dd}、\texttt{td}、\texttt{qd}、\texttt{mpf} の各版}
    \funcitem{int}{split\_dvector\_dvec\_ex}{\varname{DVector}{ret\_vec[]}, \varname{long int}{shift[]}, \varname{int}{num\_div}, \varname{DVector}{org\_vec}} $\cdots$ \bi{and the \texttt{dd}, \texttt{td}, \texttt{qd}, \texttt{mpf} counterparts; a vector carries one exponent per slice}{および \texttt{dd}、\texttt{td}、\texttt{qd}、\texttt{mpf} の各版。ベクトルはスライスごとに1つの指数を持ちます。}
\end{itemize}

%------------------------------------------------%
% End of Document
%------------------------------------------------%
